% Chapter 1

\chapter{Introduction}\label{Chapter1}

%----------------------------------------------------------------------------------------

\section{Background and Problem Statement}
Acquired Brain Injury (ABI) encompasses traumatic brain injuries (TBIs), strokes, encephalitis and hypoxic-ischemic injuries. Clinical impact is considerable. TBI alone contributes to millions of hospitalizations/emergency room visits annually \parencite{maas2017traumatic,brazinova2021epidemiology}; among working-age adults, ABI causes long-term disability, which can include physical, emotional, and cognitive sequelae \parencite{corrigan2010epidemiology}. Cognitive impairments can affect attention, memory, language, executive function, and other cognitive domains collectively; however, configurations vary substantially across patients.



Variability among patients poses a practical challenge for rehabilitation. Patients with identical injury severities may have very different cognitive profiles at follow-up evaluation following ABI, whereas common labels used for triage (diagnosis/severity band) aggregate distinctly different cognitive profiles and obscure clinically meaningful variations within groups \parencite{saatman2008classification}.



This dissertation addresses the issue empirically. Instead of applying a priori clinical taxonomy to the data set, I employ unsupervised clustering to derive cognitive subgroups from neuropsychological test scores and subsequently evaluate their clinical relevance. Clinically collected data pose an additional obstacle: missing values. Omission of full batteries occurs regularly in rehabilitation practice due to time restrictions, fatigue, examiner judgment, and patient limitations; therefore, nonrandom omission of items is expected.



Missing value imputation allows retention of incomplete assessments; however, each method introduces its own assumptions about unknown values. Limited documentation exists regarding how those assumptions will impact subsequent clustering; this dissertation evaluates that dependency directly.



Phenotyping refers to identification of cognitive subgroups via a data-driven approach. Whether those subgroups describe qualitatively distinct profiles or gradations of overall severity remains an empirical determination.

\section{Research Questions}
There are four questions I seek to answer in this dissertation:



Question RQ1: Can unsupervised clustering of the imputed data indicate discrete cognitive phenotypes in patients with acquired brain injuries?



Question RQ2: What degree does choice of imputation strategy influence emergent phenotypes?



Question RQ3: Is cluster membership predictive of clinical unit assessment? If so, this would indicate that phenotypes represent actual differences in population serviced by different units.



Question RQ4: Do phenotypes exhibit greater consistency and/or ease-of-interpretation when clustered using domain-level composite versus individual test variables?



The ultimate goal is to construct a replicable phenotyping pipeline for routinely collected clinical data. Missingness will be modeled as part of the rehabilitation process itself rather than considered a posteriori nuisance.

\section{Hypotheses}
This dissertation examines the following four hypotheses:



Hypothesis H1: Using a UMAP/HDBSCAN pipeline, I anticipate finding at least two separate stable cognitive clusters for most imputation methods, with a noise point rate below 30\% and a silhouette score above 0.40.



Hypothesis H2: Agreement should exist between cluster solutions generated using different imputation techniques (ARI/NMI \(> 0.50\)). High levels of agreement would confirm that the phenotypes do not arise solely from artifacts generated by the imputation technique(s).



Hypothesis H3: Chi-square tests and Cramér's V effect sizes (\(> 0.10\)) should illustrate a statistically significant relationship between cluster membership and clinical unit, indicating that cluster membership reflects actual differences in populations serviced by different clinics.



Hypothesis H4: Aggregating at the domain level should produce representations that are more interpretable than aggregating at the variable level, thus facilitating comparable silhouette scores and decreased multicollinearity represented by reduced condition numbers and variance inflation factors (mean VIF 2.77 to 1.88; mean CN 61.7 to 11.8).



CogDash is an interactive dashboard designed to execute this entire pipeline for future patients; additionally it provides both severity tier assignment and consensus confidence across all imputation techniques. The code/notebooks/dashboard are located at \url{https://github.com/zkunos/thesis_project}.
